\vspace{-2mm}\section{Introduction}\vspace{-1mm}

\begin{figure}[t]
    \centering
    \includegraphics[width=0.98\textwidth]{figures/framework.pdf}
    \caption{Overview of our study: Automatic Promotion (AutoPR) task, its benchmark PRBench, and the associated method PRAgent. The details of citation trend analysis are shown in Appendix~\ref{append:citation}.
    }
    \label{fig:framework}
\end{figure}

The rapid development of Large Language Models (LLMs) has pushed the capabilities of artificial intelligence to new heights~\cite{achiam2023gpt,yang2025qwen3,qin2024large}.
In particular, LLMs that enhance complex task-solving abilities through specific training strategies, such as Chain-of-Thought (CoT), are at the forefront of this progress~\cite{wei2022chain,guo2025deepseek,jaech2024openai}.
We refer to these as \textbf{Large Reasoning Models (LRMs)}~\cite{chen2025towards}.
However, this advancement is accompanied by a growing and significant challenge: Hallucination.

Initially, research on hallucination was primarily concentrated on the factual level~\cite{huang2025survey,park2025steer}.
For instance, early surveys classified hallucinations into two main categories:
Factuality Hallucination, where generated content conflicts with verifiable real-world knowledge, and Faithfulness Hallucination, where the output contradicts the provided source text or context.
Although significant, these errors are mainly concerned with the correctness of the output's result.

However, the nature of hallucination has fundamentally shifted as LRMs increasingly perform multi-step reasoning, particularly driven by techniques like CoT.
The CoT strategy was introduced to improve complex problem-solving by guiding models to generate explicit intermediate steps.
While this approach appears to offer a path to more reliable and interpretable AI, deeper investigation reveals a concerning reality.
The explicit reasoning paths, while offering transparency, simultaneously introduce new and subtle vulnerabilities.
The core of the hallucination problem has shifted from \textit{what the model knows} to \textit{how the model thinks}.
This has led to the emergence of a more complex and concealed category: \textbf{Reasoning Hallucination}.

In this paradigm, a model may produce the correct final answer, yet its reasoning process is flawed with logical fallacies, fabrications, or disconnected steps.
Conversely, a model might follow a seemingly coherent reasoning chain only to arrive at an incorrect conclusion~\cite{joshi2024llms,yao2025reasoning}.
Specifically, the reasoning process of these models can be unreliable.
This unreliability manifests not only as logical fallacies but also as a pathological disregard for simple solutions, where models fall into unnecessary complexity and overthinking~\cite{peng2025revisiting,fan2025missing}.
This behavior is perhaps another byproduct of their reasoning-style text pattern matching mechanism.
This transition from focusing on factual errors to reasoning process flaws represents the current research frontier~\cite{jiao2025trustworthy}.
It also highlights the profound paradoxes underlying advanced reasoning capabilities.

Ironically, while CoT can reduce the frequency of factual hallucinations in some cases, it may also obscure key cues for their detection.
When the model generates a long and complex reasoning chain, its reported confidence in the final output increases significantly, and the extended reasoning process itself appears more persuasive.
This makes hallucinations harder to detect, both for automated evaluation toolkits and human reviewers~\cite{cheng2025chain}.
This creates a fundamental tension between capability and evaluability, which we term the \textbf{Reasoning Evaluability Paradox}.
This paradox is that the very tools we use to enhance model reasoning (like CoT) may simultaneously weaken our ability to assess their reliability.

The main contributions of this work are as follows:

\begin{itemize}[leftmargin=16pt, itemsep=0pt, topsep=0pt]
    \item \textbf{A Systematic Taxonomy of Reasoning Hallucination:}
    We propose a systematic taxonomy for \textbf{Reasoning Hallucination}.
    This framework moves beyond traditional, result-based definitions of hallucination to focus on process-based flaws that occur during complex reasoning.
    We subdivide reasoning hallucination into four main categories (\textbf{Logical Fallacies, Flaw Repetition, Think-Answer Mismatch, and Overthinking}) to facilitate detailed analysis.

    \item \textbf{A Comprehensive Review of the Reasoning Paradigm:}
    This paper provides a comprehensive review of the causes, detection methods, and mitigation strategies for reasoning hallucination.
    We place special emphasis on analyzing how techniques like CoT fundamentally alter the form of hallucination, shifting it from incorrect answers to flawed reasoning processes.

    \item \textbf{Critical Analysis and Future Roadmap:}
    We conduct a critical analysis of the core challenges in current research, particularly the \textbf{Reasoning Evaluability Paradox} as defined in this paper.
    Finally, we propose future research directions for building more robust, reliable, and trustworthy large reasoning models.
\end{itemize}
